\documentclass[12pt]{amsart}  
 
\usepackage{amsmath,amssymb,amsfonts,amscd,xcolor,tikz-cd}
\usepackage[hidelinks]{hyperref}
\numberwithin{equation}{section}

\definecolor{commentred}{RGB}{190,0,0}
\definecolor{additionpurple}{RGB}{120,0,160}
\newcommand{\redcomment}[1]{\par\smallskip\noindent{\color{commentred}\textbf{Comment.} #1}\par\smallskip}
\newcommand{\redcorr}[1]{{\color{commentred}#1}}
\newcommand{\purpledetail}[1]{{\color{additionpurple}#1}}
\newcommand{\purplecomment}[1]{\par\smallskip\noindent{\color{additionpurple}\textbf{Added detail.} #1}\par\smallskip}
\newenvironment{redadd}{\par\smallskip\begingroup\color{commentred}}{\par\smallskip\endgroup}
\newenvironment{purpleadd}{\par\smallskip\begingroup\color{additionpurple}}{\par\smallskip\endgroup}

\providecommand{\k}{}\renewcommand{\k}{\mathbf{k}}
\providecommand{\C}{}\renewcommand{\C}{\mathcal{C}}
\providecommand{\D}{}\renewcommand{\D}{\mathcal{D}}
\providecommand{\M}{}\renewcommand{\M}{\mathcal{M}}
\providecommand{\Vec}{}\renewcommand{\Vec}{\mathrm{Vec}}
\providecommand{\sVec}{}\renewcommand{\sVec}{\mathrm{sVec}}
\providecommand{\Ver}{}\renewcommand{\Ver}{\mathrm{Ver}}
\providecommand{\Cl}{}\renewcommand{\Cl}{\mathrm{Cl}}
\providecommand{\h}{}\renewcommand{\h}{\mathfrak{h}}
\providecommand{\ot}{}\renewcommand{\ot}{\otimes}
\providecommand{\gr}{}\renewcommand{\gr}{\operatorname{gr}}
\providecommand{\id}{}\renewcommand{\id}{\operatorname{id}}
\providecommand{\Hom}{}\renewcommand{\Hom}{\operatorname{Hom}}
\providecommand{\End}{}\renewcommand{\End}{\operatorname{End}}
\providecommand{\Pic}{}\renewcommand{\Pic}{\operatorname{Pic}}
\providecommand{\Br}{}\renewcommand{\Br}{\operatorname{Br}}

\theoremstyle{plain}
\newtheorem{theorem}{Theorem}[section]
\newtheorem{lemma}[theorem]{Lemma}
\newtheorem{proposition}[theorem]{Proposition}
\newtheorem{corollary}[theorem]{Corollary}

\theoremstyle{definition}
\newtheorem{definition}[theorem]{Definition}
\newtheorem{example}[theorem]{Example}
\newtheorem{question}[theorem]{Question}
\newtheorem{conjecture}[theorem]{Conjecture}

\theoremstyle{remark}
\newtheorem{remark}[theorem]{Remark}

\pagestyle{plain}

\begin{document}

\title{Clifford and Weyl algebras in symmetric tensor categories}

\author{Pavel Etingof}
\address{Department of Mathematics,
Massachusetts Institute of Technology,
77 Massachusetts Avenue,
Cambridge, MA 02139,
USA
}
\email{etingof@math.mit.edu}


\begin{abstract} Let $\C$ be a symmetric tensor category over an algebraically closed field $\k$ of characteristic $\ne 2$. We study Clifford and Weyl algebras of objects of $\C$ with a (skew-)symmetric bilinear form. When the form is non-degenerate, we establish simplicity and the Azumaya property for such algebras under suitable assumptions. We also compute Clifford and Weyl algebras in the Verlinde category $\Ver_p$ and use them to prove that if $\C$ is Frobenius exact then the Weyl algebra of a symplectic object of $\C$ with finite symmetric algebra is Azumaya.  Using this, we introduce the symplectic Witt group $\mathcal S\mathcal W(\C)$, the subgroup of the Brauer group $\Br(\C)$ consisting of Morita classes of such Azumaya algebras, and when $\C={\rm Rep}(G)\boxtimes\sVec$ for a finite group $G$ of order coprime to ${\rm char}(\k)$, express $\mathcal S\mathcal W(\C)$ in terms of second Stiefel-Whitney classes of orthogonal representations of $G$.
\end{abstract}


\maketitle

\tableofcontents

\section{Introduction} 

Clifford and Weyl algebras, which are quantizations of exterior and symmetric algebras respectively, are among the simplest and most important examples of non-commutative algebras. In this paper we develop a theory of Clifford and Weyl algebras in an arbitrary symmetric tensor category $\C$ and establish their basic properties (such as the  correspondence between Clifford and Weyl algebras by parity change and the PBW theorem). We also establish the simplicity of these algebras under suitable assumptions; in particular, we show that Weyl algebras of objects with finite symmetric algebras in Frobenius exact categories are Azumaya algebras. Finally, we define the symplectic Witt group $\mathcal S\mathcal W(\C)$ of $\C$ - the subgroup of the Brauer group $\Br(\C)$ consisting of Morita classes of such Azumaya algebras, 
and describe it explicitly in the case of the category of super-representations of a finite group of order coprime to ${\rm char}(\k)$ in terms of second Stiefel-Whitney classes of its orthogonal representations. 

The organization of the paper is as follows. In Section 2 we discuss the theory of Clifford and Weyl algebras of arbitrary (skew)-symmetric forms. In Section 3 we discuss Clifford and Weyl algebras of non-degenerate forms, define symplectic Witt groups, and compute the symplectic Witt group 
of the category of super-representations of a finite group. 

{\bf Acknowledgements.} In this paper I collaborated with ChatGPT 5.5 and 5.6 Pro. This work was partially supported by the NSF grant DMS-2502467. I am grateful to V. Ostrik for useful comments. 



\section{Clifford and Weyl algebras} 

\subsection{Definition and basic properties of Clifford and Weyl algebras} 
Throughout the paper, $\k$ will denote an algebraically closed field of characteristic $\ne 2$. Let $\C$ be a symmetric tensor category over $\k$ (\cite{EGNO}, Definitions 4.1.1, 8.1.1, 8.1.12).   
Let $V$ be an object of $\C$ equipped with a bilinear form $B\in \Hom(V\otimes V,\mathbf{1})$. Suppose $B$ is either symmetric, i.e., $B\circ c_{V,V}=B$, or skew-symmetric, i.e., $B\circ c_{V,V}=-B$, where $c$ is the symmetric braiding of $\C$. In this case $\h_V:=V\oplus \mathbf{1}$ has a natural structure of a Lie superalgebra in $\C$, i.e., 
the image of a Lie algebra in $\C\boxtimes \sVec$ under the (non-symmetric) forgetful functor to $\C$. Namely,
if $B$ is skew-symmetric then $\h_V$ is purely even and is called the {\it Heisenberg algebra} of $V$; 
the commutator maps $V\otimes V$ to $\mathbf{1}$ by $B$, and $\mathbf{1}$ is central. If $B$ is symmetric, 
the definition is the same, except $\mathbf{1}$ is even but $V$ is odd; in this case we say that 
$\h_V$ is the {\it Heisenberg superalgebra} of $V$. 

Define the associative (ind-)algebra $A(V)$ in $\C$ 
to be the quotient of the enveloping algebra $U(\h_V)$ by the relation identifying the unit 
copy of $\mathbf{1}$ in $U(\h_V)$ with $\mathbf{1}\subset \h_V$. In other words, if $z_U:\mathbf 1\to\mathfrak h_V\to U(\mathfrak h_V)$ is the central generator and $1_U:\mathbf 1\to U(\mathfrak h_V)$ is the algebra unit, then 
$$
A(V):=U(\mathfrak h_V)/({\rm Im}(z_U-1_U)).
$$
We call 
$A(V)$ the {\it Weyl algebra} of $V$ if $B$ is skew-symmetric and 
the {\it Clifford algebra} of $V$ if $B$ is symmetric. 
For $\C=\Vec$, this agrees with the definitions 
in classical algebra.\footnote{Note that many classical references normalize the Clifford algebra relation as $vw+wv=2B(v,w)$. In this paper we drop the factor of $2$.} 

The relations of $A(V)$ are invariant under the automorphism $-{\rm Id}$ on $V$, hence 
$A(V)=A_+(V)\oplus A_-(V)$, where this automorphism acts on $A_\pm(V)$ by $\pm 1$. 
In the Clifford algebra case, it is natural to think of this $\Bbb Z/2$-grading as a super-grading. 
Note also that changing the sign of $B$ transforms $A(V)$ into $A(V)^{\rm op}$, where in the Clifford case ${\rm op}$ 
denotes the super-opposite algebra. Thus for every choice of $i=\sqrt{-1}\in \k$ we obtain an isomorphism 
$A(V)\cong A(V)^{\rm op}$ (as rescaling each argument by $i$ transforms $B$ to $-B$). 

Let $\psi\in \C$ be a super-line, i.e., an invertible object of dimension $-1$ with an identification $\alpha: \psi\otimes \psi\cong \mathbf{1}$ (so $c_{\psi,\psi}=-\id_{\psi\otimes\psi}$); for example,
there is a natural one if $\C=\C_+\boxtimes \sVec$, where $\C_+$ is another tensor category and $\psi$ is the generating simple object of $\sVec$. We will identify $\psi\otimes \psi$ with ${\mathbf 1}$ using $\alpha$. 
If an object $V$ is equipped with a bilinear form $B$, then so is $V\otimes \psi$, but this 
functor exchanges symmetric and skew-symmetric forms. 

Let 
$$
A_\psi(V):=A_+(V)\oplus A_-(V)\otimes \psi
$$
with the natural multiplication.
Explicitly, the product on $A_\psi(V)$ restricts to the original product on $A_+(V)$ and the $A_+(V)$-bimodule structure on $A_-(V)$, while the product of two odd terms is
\[
(A_-(V)\otimes\psi)\otimes(A_-(V)\otimes\psi)\xrightarrow{\;(m\otimes\alpha)(\id\otimes c_{\psi,A_-(V)}\otimes \id)\;}A_+(V).
\]


\begin{lemma}\label{leweyl1}  (i) There is a natural isomorphism 
$
A(V\otimes \psi)\cong A_\psi(V).
$

(ii) Let $V,W\in \C$ be two objects equipped with bilinear forms of the same type, i.e., both symmetric or both skew-symmetric. Then 
$$
A(V)\otimes A(W)\cong A(V\oplus W),
$$ 
where in the case of symmetric forms this is the super-tensor product.
\end{lemma} 

\begin{proof}
For (i), define the form on $V\otimes\psi$ by
\[
(V\otimes\psi)^{\otimes2}\xrightarrow{\;\id_V\otimes c_{\psi,V}\otimes\id_\psi\;}V^{\otimes2}\otimes\psi^{\otimes2}
\xrightarrow{\;B\otimes\alpha\;}\mathbf{1}.
\]
Because $c_{\psi,\psi}=-\id$, this changes the symmetry type of $B$. The generating object $V\otimes\psi$ is sent to the odd part $A_-(V)\otimes\psi$ of $A_\psi(V)$, and the defining commutation relation is exactly as above. Hence the universal property of the enveloping-algebra quotient gives $A(V\otimes\psi)\cong A_\psi(V)$.

 (ii) In the Weyl case, the generating objects $V$ and $W$ commute because the cross-pairing is zero, so the universal property gives 
$$
A(V)\otimes A(W)\cong A(V\oplus W).
$$ 
In the Clifford case the same argument is made in superalgebras: the cross-pairing is zero, so the odd generating objects from the two summands super-commute, which is precisely the super-tensor product convention.
\end{proof} 

Lemma \ref{leweyl1}(i) shows that in symmetric tensor categories containing a super-line, the theories of Weyl algebras and Clifford algebras are equivalent. 

\subsection{Contraction maps} 

Recall that for any $W\in \C$, the symmetric algebra $SW$ is a Hopf algebra with coproduct 
decomposing into homogeneous components $\Delta_{i,j}: S^{i+j}W\to S^iW\otimes S^jW$. 
In particular, by duality the components $\Delta_{1,j-1}$ define maps 
$$
D_j: W^*\otimes S^jW\to S^{j-1}W
$$ 
which combine into the contraction map 
$$
D: W^*\otimes SW\to SW.
$$ 
Thus if $V\in \C$ is equipped with a bilinear form $B$ then we have a map 
$$
D_{B,ij}: S^iV\otimes S^jV\to S^{i-1}V\otimes S^{j-1}V
$$
defined by 
$$
D_{B,ij}:=(D_i\otimes D_j)(\id\otimes c_{V^*,S^iV}\otimes \id)(B^*\otimes \id\otimes \id),
$$
where $B^*: \mathbf{1}\to V^*\otimes V^*$ is dual to $B$. 
These maps 
combine into the bicontraction map
$$
D_B: SV\otimes SV\to SV\otimes SV.
$$  

More generally, given $k\in \mathbb Z_{\ge 0}$, subsets $S\subset [1,i],T\subset [1,j]$ with 
$|S|=|T|=k$, and a bijection $\sigma: S\to T$, we have 
a morphism 
$$
B^{S,T,\sigma}_{ijk}: V^{\otimes i}\otimes V^{\otimes j}\to S^{i-k}V\otimes S^{j-k}V
$$
obtained by contracting the $S$-tensorands of $V^{\otimes i}$ with $T$-tensorands 
of $V^{\otimes j}$ according to $\sigma$ using $B$ and then projecting to symmetric powers. 
Then the morphism 
$$
B_{ijk}:=\sum_{S,T,\sigma} B^{S,T,\sigma}_{ijk}
$$
is $\mathfrak S_i\times \mathfrak S_j$-invariant, hence descends to a morphism 
$$
D_{B,ij}^{(k)}: S^iV\otimes S^jV\to S^{i-k}V\otimes S^{j-k}V.
$$ 
These morphisms combine into the map 
$$
D_B^{(k)}: SV\otimes SV\to SV\otimes SV.
$$ 
It is easy to see that 
$$
D_B^k=k! D_B^{(k)},
$$
in particular if ${\rm char}(\k)=p>0$ then $D_B^p=0$. 

\subsection{The PBW theorem}
The algebra $A(V)$ has an increasing $\mathbb Z_{\ge 0}$-filtration $F^\bullet$ defined by $\deg(V)=1$, and we have a natural surjective  homomorphism of graded algebras
$\xi: SV\to {\rm gr}A(V)$ for Weyl algebras and $\xi: \wedge V\to {\rm gr}A(V)$ for Clifford algebras. 

\begin{proposition}\label{pbw} (PBW theorem) $\xi$ is an isomorphism. Thus $A(V)$ is a filtered quantization of the  Poisson 
algebra $SV$ or super-Poisson algebra $\wedge V$. 
\end{proposition} 

\begin{proof} By replacing $\C$ with $\C\boxtimes \sVec$ if necessary, we may assume without loss of generality that $\C$ contains a super-line. Then by Lemma \ref{leweyl1}(i), it suffices to prove the statement for Weyl algebras.
In this case, there is a natural morphism
$$
\beta_k: SV\otimes SV\to SV
$$ 
of degree $-2k$ given by 
$$
\beta_k:=m_0\circ D_B^{(k)},
$$
where $m_0$ is the multiplication on $SV$. For example, $\beta_0$ is the usual product and $\beta_1$ is the Poisson bracket defined by $B$. 

We can now define the {\it Moyal-Weyl} product on $SV$ by the formula 
$$
m:=\sum_{k\ge 0}2^{-k}\beta_k=m_0\circ E,\ E:=\sum_{k\ge 0}2^{-k}D_B^{(k)}.
$$
(Note that in characteristic zero $E$ can be expressed by the usual Moyal-Weyl formula
$E=\exp(\tfrac{1}{2}D_B)$).
It is easy to see that 
$$
E\circ (m_0\otimes \id)=(m_0\otimes \id) \circ E_{13}E_{23}
$$
as maps $SV\otimes SV\otimes SV\to SV\otimes SV$, where $E_{rs}$ denotes $E$ applied to the $r$-th and $s$-th factors. Namely, this identity follows from the fact that contraction with a tensor product is the sum of contractions with its factors. 

It follows that the product $m$ is associative. Indeed, 
$$
m\circ (m\otimes \id)=m_0\circ E\circ (m_0\otimes \id)\circ (E\otimes \id)=
m_0\circ (m_0\otimes \id) \circ E_{13}E_{23}E_{12}
$$
and likewise 
$$
m\circ (\id\otimes m)=m_0\circ (\id \otimes m_0) \circ E_{12}E_{13}E_{23}.
$$
Thus the statement follows from associativity of $m_0$ and the fact that $E_{ij}$ commute.

It follows that there is a natural homomorphism $\phi: A(V)\to (SV,m)$ sending $V$ to $V$ by the identity map, and we have $\operatorname{gr}(\phi)\circ \xi=\id_{SV}$. So $\xi$ is injective, hence an isomorphism. 
\end{proof} 

\begin{remark} Here is a shorter but less explicit proof of Proposition \ref{pbw} after reduction to skew-symmetric forms. Given an object $V\in \C$ with a skew-symmetric form $B$, define the affine group scheme $\mathcal H_V$ in $\C$, which 
as a scheme is the object $V\oplus {\mathbf 1}$ (i.e., $\mathcal O(\mathcal H_V)=S(V^*\oplus {\mathbf 1})$) with group law given by the following schematic formula (i.e., formula at the level of points over commutative algebras in $\C$): 
$$
(v_1,\alpha_1)*(v_2,\alpha_2):=(v_1+v_2,\alpha_1+\alpha_2+\tfrac{1}{2}B(v_1,v_2)).
$$
We call this group the {\it Heisenberg group} of $V$. Then $\h_V={\rm Lie}\mathcal H_V$. Thus 
the proposition follows from
\cite{E}, Example 4.8(2). 
\end{remark} 

\section{Clifford and Weyl algebras of non-degenerate forms} 
From now on assume that $B$ is non-degenerate, i.e., defines an isomorphism $V\to V^*$ by evaluating the first (or, equivalently, second) factor.\footnote{Indeed, the two morphisms
$B_1,B_2:V\longrightarrow V^*$ obtained by evaluating the first and the second tensorand of $B$ respectively are transposes of one another under rigidity (with the symmetry inserted). Hence $B_1$ is an isomorphism if and only if $B_2$ is an isomorphism.}

\subsection{Simplicity of Clifford and Weyl algebras} 

\begin{proposition}\label{simplic} Let $R$ be a simple (ind-)algebra in $\C$.

(i) Suppose $V$ is equipped with a skew-symmetric bilinear form $B$, and $S^p V=0$  (hence $S^iV=0$ for $i\ge p$) if $p={\rm char}(\k)>0$.  Then the algebra $A_R(V):=R\otimes A(V)$ is simple.

(ii) Suppose $V=W\oplus W^*$ with the skew-symmetric form defined by the pairing between $W$ and $W^*$, and 
$S^pW=S^pW^*=0$ (hence $S^iW=S^iW^*=0$ for $i\ge p$) if $p={\rm char}(\k)>0$.  Then the algebra $A_R(V):=R\otimes A(V)$ is simple.
\end{proposition}

\begin{example} Let $\C=\sVec$, $V\in \C$ a purely odd space (i.e., $V=V_0\otimes \psi$, where $V_0$ is a usual vector space over $\k$), and $R=\k$. Then 
Proposition \ref{simplic}(i) states that for a non-degenerate symmetric form $B_0$ on $V_0$, the Clifford algebra ${\rm Cl}(V_0,B_0)$ is simple {\it as a super-algebra}. Note that if $\dim V_0$ is odd, 
it is not simple as an ordinary algebra; rather, it is the direct sum of two matrix blocks. 
\end{example} 

\begin{proof} The filtration on $A(V)$ extends naturally to $A_R(V)$ by setting $\deg R=0$. 

(i) Let $i>0$, and identify $V$ with $V^*$ by evaluating $B$ on the first argument. Recall that $D_i: V\otimes S^iV\to S^{i-1}V$ corresponds by duality to the map $\Delta_{1,i-1}: S^iV\to V\otimes S^{i-1}V$.
Composing $\Delta_{1,i-1}$ with multiplication $V\otimes S^{i-1}V\to S^iV$ gives multiplication by $i$. Hence  in characteristic zero or for $i<p$ in characteristic $p$ the map $\Delta_{1,i-1}$ is a split injection. 

It follows that for such $i$, if $E\subset R\otimes S^iV$ is a nonzero object, then $(1_R\otimes D_i)(V\otimes E)\ne 0$.
Thus if $E\subset F^iA_R(V)$ is a nonzero subobject whose image in
$\operatorname{gr}^iA_R(V)=F^iA_R(V)/F^{i-1}A_R(V)$
is nonzero, then the leading term of $[\,,\,](V\otimes E)$ is the contraction above, so the subobject $[\,,\,](V\otimes E)\subset F^{i-1}A_R(V)$ is nonzero. Now let $I$ be a nonzero two-sided ideal in $A_R(V)$ and suppose $i$ is the smallest integer such that $I\cap F^iA_R(V)\ne 0$. Denote this intersection by $E$. If $i>0$ then as we have shown, $[\,,\,](V\otimes E)\subset I\cap F^{i-1}A_R(V)$ is nonzero, a contradiction with minimality of $i$. Thus $i=0$, so $I\cap R\ne 0$. Since $R$ is simple, $I\cap R=R$, so $I$ contains the unit of $A_R(V)$, hence $I=A_R(V)$. 

(ii) By Proposition \ref{pbw}, the ordered multiplication map
\[
R\otimes SW\otimes SW^*\longrightarrow A_R(W\oplus W^*)
\]
is an isomorphism of objects: its associated graded map is the identity on $R\otimes SW\otimes SW^*$. Define a filtration of $F^dA_R(W\oplus W^*)$ by \scriptsize
\[ 
F^{d,b}A_R(W\oplus W^*):=R\otimes\left(\bigoplus_{a+c<d}S^aW\otimes S^cW^*
\oplus\bigoplus_{c\le b}S^{d-c}W\otimes S^cW^*\right),
\ 0\le b\le d,
\] \normalsize
where the right-hand side is transported to $A_R(W\oplus W^*)$ by ordered multiplication.
The successive subquotients are
\[
\operatorname{gr}^{d,b}A_R(W\oplus W^*)=F^{d,b}/F^{d,b-1}\cong R\otimes S^{d-b}W\otimes S^bW^*,
\]
where $F^{d,-1}:=R\otimes\bigoplus_{a+c<d}S^aW\otimes S^cW^*$ (and $F^{0,-1}:=0$). Let $I$ be a nonzero two-sided ideal in $A_R(W\oplus W^*)$, and let $(d,b)$ be minimal in lexicographic order such that $I\cap F^{d,b}\ne0$. Let $E$ be the nonzero image of $I\cap F^{d,b}$ in $\operatorname{gr}^{d,b}A_R(W\oplus W^*)$.

If $b>0$ then the contraction map 
$$
D_{b,W^*}: W\otimes S^{d-b}W\otimes S^bW^*\to S^{d-b}W\otimes S^{b-1}W^*
$$ 
corresponds by duality to the injective map 
$$
\id\otimes \Delta_{1,b-1}^{W^*}: S^{d-b}W\otimes S^bW^*\to S^{d-b}W\otimes W^*\otimes S^{b-1}W^*
$$ 
(namely, in characteristic $p$ since $E\ne0$, we have $0<b<p$, and $\Delta^{W^*}_{1,b-1}$ followed by multiplication is $b\cdot \id$, so it is injective). It follows that  $(1_R\otimes D_{b,W^*})(W\otimes E)\ne 0$.
The leading term of the commutator of $W$ with $F^{d,b}$ is this contraction, so $[W,I]$ has a nonzero image in total degree $d-1$, contradicting minimality of $d$.

Similarly, if $b<d$ then the contraction map  
$$
D_{d-b,W}: W^*\otimes S^{d-b}W\otimes S^bW^*\to S^{d-b-1}W\otimes S^bW^*
$$ 
corresponds by duality to the injective map 
$$
\Delta_{d-b-1,1}^W\otimes \id: S^{d-b}W\otimes S^bW^*\to S^{d-b-1}W\otimes W\otimes S^bW^*.
$$ 
So $(1_R\otimes D_{d-b,W})(W^*\otimes E)\ne 0$. 
The leading term of the commutator of $W^*$ with $F^{d,b}$ is this contraction, so $[W^*,I]$ has a nonzero image in total degree $d-1$, again a contradiction. Thus $d=0$. Hence $I\cap R\ne 0$, and since $R$ is simple, $I\cap R=R$; consequently $I$ contains the unit, hence $I=A_R(W\oplus W^*)$.
\end{proof} 

\begin{corollary}\label{simplic1} Let $\operatorname{char}(\k)=p>0$ and 
$$
V=V_1\oplus\cdots\oplus V_n,\ B=B_1\oplus\cdots\oplus B_n,
$$ 
where $B_i$ are skew-symmetric bilinear forms on $V_i$. If for all $1\le m\le n$ either $S^pV_m=0$ or $V_m=W_m\oplus W_m^*$ with $B_m$ defined by the pairing of $W_m,W_m^*$ with $S^pW_m=S^pW_m^*=0$ then $A(V)$ is a simple algebra. 
\end{corollary} 

\begin{proof} The proof is by induction in $n$ starting with the trivial case $n=0$. Take $n>0$. Let $V'=V_1\oplus\cdots\oplus V_{n-1}$. Then by Lemma \ref{leweyl1}(ii), $A(V)=A(V')\otimes A(V_n)$. But $A(V')$ is simple by the induction assumption, so $A(V)$ is simple by Proposition \ref{simplic}. 
\end{proof} 

Observe that the algebra $A(W\oplus W^*)$ acts naturally on $SW$, with $W$ acting by multiplication and 
$W^*$ by contraction. Thus we have a homomorphism 
$$
\eta: A(W\oplus W^*)\to \underline{\End}(SW).
$$ 

Also recall from \cite{OZ} that an {\it Azumaya algebra} in $\C$ is an algebra $A$ such that 
the natural map $A\otimes A^{\rm op}\to \underline{\End}(A)$ is an isomorphism. Morita equivalence classes of 
Azumaya algebras form a group called the {\it Brauer group} of $\C$, denoted $\Br(\C)$. By \cite{DN}, Subsection 5.5, 
if $\C$ is finite then there is a natural isomorphism $\Br(\C)\to \Pic(\C)$ which sends $A$ to $A$-mod. 

\begin{corollary}\label{simplic2} Let $\C$ be a Frobenius exact symmetric tensor category (cf. \cite{CEO}).

(i) If $V\in \C$ is equipped with a skew-symmetric form and the symmetric algebra $SV$ has finite length then $A(V)$ is a simple algebra.

(ii) If $W\in \C$ and the symmetric algebra $SW$ has finite length then $\eta$ is an isomorphism.

(iii) In (i) $A(V)$ is a (central) Azumaya algebra. 
\end{corollary}

\begin{proof} Recall that a symmetric tensor category $\D$ is said to have {\it moderate growth} if for every 
$X\in \D$ the length $\ell(X^{\otimes n})$ grows at most exponentially with $n$. 
We will need the following lemma. 

\begin{lemma}\label{lemo} Let $\D$ be a symmetric tensor category over $\k$ tensor generated by an object $V$ (i.e., 
every object of $\D$ is a subquotient of a finite direct sum of tensor words in $V$ and $V^*$). If the algebra 
$SV$ or $\wedge V$ has finite length then $\D$ has moderate growth. 
\end{lemma} 

\begin{proof} We may assume that $SV$ has finite length: otherwise, if $\wedge V$ has finite length, 
replace $\D$ with $\D\boxtimes \sVec$ and $V$ by $\psi\oplus V\otimes \psi$ (the category tensor generated by this object is $\D\boxtimes \sVec$). Thus $S^NV=0$ for some $N$. So in the case ${\rm char}(\k)=0$ the category $\D$ 
is Schur-finite, hence has moderate growth (\cite{De}); in fact, by the main result of 
\cite{De} it is super-Tannakian. 

So we may assume that ${\rm char}(\k)=p>0$. 
Then the symmetrizing element
\[
\mathbf e_N:=\sum_{\sigma\in\mathfrak S_N}\sigma\in\k\mathfrak S_N
\]
factors through $V^{\otimes N}\twoheadrightarrow S^NV$, and hence acts by zero on $V^{\otimes N}$. Since $\mathbf e_N\ne0$ in $\k\mathfrak S_N$, this gives $\operatorname{sd}(V)<\infty$ in the notation of \cite[Definition~4.1]{CEO}.
Since duality intertwines the symmetric-group actions on $V^{\otimes n}$ and $(V^*)^{\otimes n}$ (up to the involution $\sigma\mapsto\sigma^{-1}$ of $\k\mathfrak S_n$), it follows that $\operatorname{sd}(V^*)=\operatorname{sd}(V)<\infty$. Hence \cite[Proposition~4.7(6)]{CEO} gives $\operatorname{ad}(V),\operatorname{ad}(V^*)<\infty$. Lemma~4.9(2) of \cite{CEO} then implies that
$Z:=\mathbf1\oplus V\oplus V^*$ has finite alternating degree, and hence finite growth degree by \cite[Proposition~4.7(5)]{CEO}. For any fixed $X\in\D$, there are $r,M\ge1$ such that $X$ is a subquotient of $(Z^{\otimes r})^{\oplus M}$. Therefore $X^{\otimes n}$ is a subquotient of $(Z^{\otimes rn})^{\oplus M^n}$, and
\[
 \ell(X^{\otimes n})\le M^n\ell(Z^{\otimes rn}),
\]
which grows at most exponentially. Thus $\D$ has moderate growth.
\end{proof} 

Now we proceed to prove the corollary. Without loss of generality, we may assume that $\C$ is tensor generated by $V$ in (i),(iii) and by $W$ in (ii): otherwise replace $\C$ by the full tensor subcategory generated by $V$ or $W$. Then by Lemma \ref{lemo}, $\C$ has moderate growth. Hence if ${\rm char}(\k)=0$ then $\C$ is super-Tannakian by \cite{De}, thus all parts of the corollary follow from classical algebra. So we may assume that ${\rm char}(\k)=p>0$. 
In this case, by \cite{CEO}, $\C$ admits a fiber functor $F: \C\to \Ver_p$.

(i) Since $F$ is exact and faithful, it suffices to show that $F(A(V))=A(F(V))$ is simple. Now, $F(V)=\oplus_{i=1}^{p-1}M_i\otimes L_i$, where $M_i$ are finite-dimensional $\k$-vector spaces and $L_i$ are simple objects of $\Ver_p$. Since $SV$ has finite length, so does $SF(V)=F(SV)$, so $M_1=0$.
But for $i>1$, $S^pL_i=0$ (\cite[Section~6]{EOV}). Also, $L_i$ is self-dual, and its invariant bilinear form is skew-symmetric for $i$ even and symmetric for $i$ odd. Thus a skew form on the isotypic component $M_i\otimes L_i$ corresponds to a symmetric form on $M_i$ if $i$ is even, and to a skew form on $M_i$ if $i$ is odd. Hence the form $B$ decomposes into a direct sum of forms on $V_i=L_i$ with $i$ even and forms on $V_j=L_j\oplus L_j$ for odd $j$ pairing the factors. More explicitly, for $i$ even the multiplicity space $M_i$ carries a non-degenerate symmetric form, so it is an orthogonal sum of lines, giving copies of $L_i$. For $i$ odd the multiplicity space $M_i$ carries a non-degenerate skew-symmetric form, so it is a sum of symplectic planes, giving summands $L_i\oplus L_i$. Thus the simplicity of $A(V)$ follows from Corollary \ref{simplic1}. 

(ii) It suffices to prove the statement after applying $F$, i.e., in $\Ver_p$. We have $F(W)=\oplus_{i=1}^{p-1}N_i\otimes L_i$, and since $SW$ has finite length, $N_1=0$. Thus by Lemma \ref{leweyl1}(ii) it suffices to prove the statement for 
$W=L_i$, $i\ge 2$.

But $S^pL_i=0$, hence by Proposition~\ref{simplic}(ii) $A(W\oplus W^*)$ is simple. So $\eta$ is injective. By Proposition \ref{pbw}, the multiplication map 
$$
SW^*\otimes SW\to A(W\oplus W^*)
$$ 
is an isomorphism in $\C$, and $\underline{\End}(SW)\cong (SW)^*\otimes SW=SW^*\otimes SW$, so the source and target of $\eta$ have the same length. Thus $\eta$ is an isomorphism. 

(iii) The opposite algebra to $A(V)=A(V,B)$ is $A(V,-B)$, and the orthogonal direct sum $(V,B)\oplus(V,-B)$ is isometric, using the isomorphism $V\xrightarrow{\sim}V^*$ given by $B$, to the object $V\oplus V^*$ with the skew-symmetric form defined by the pairing of the summands. Under the PBW isomorphism, the regular left-right action
\[
 A(V)\otimes A(V)^{\rm op}\longrightarrow\underline{\End}(A(V))
\]
identifies with the map
\[
 A(V\oplus V^*)\longrightarrow\underline{\End}(SV)
\]
from (ii). Hence it is an isomorphism.
\end{proof} 

\subsection{The top symmetric power of a symplectic object}

Let $\C$ be a Frobenius exact symmetric tensor category. 
Let $V\in  \C$ be a symplectic object (i.e., equipped with a non-degenerate skew-symmetric form) 
and assume that $SV$ has finite length. Following \cite{CEN}, let $m(V)$ be the largest integer $m$ for which $S^mV\ne 0$. 
By \cite[Proposition~2.3.2]{CEN}, the object $\psi=S^{m(V)}V$ is invertible.

\begin{lemma}\label{psisq} (i) For each $0\le i\le m(V)$ there are natural isomorphisms 
$$
S^iV\cong (S^iV)^*\cong S^{m(V)-i}V\otimes \psi
$$ 
and $\psi\otimes \psi\cong {\mathbf 1}$. 

(ii) $c_{\psi,\psi}=(-1)^{m(V)}\id_{\psi\otimes \psi}$. In particular, if $m(V)$ is odd, then $\psi$ is a super-line. 
\end{lemma} 
 
\begin{proof} We may assume $i\le m(V)$. 

(i) By Corollary \ref{simplic2}(ii), $SV$ is a simple module over $A(V\oplus V^*)$, hence generated by $\psi$. Thus
$$
(SV^*)\psi=(SV^*\cdot SV)\psi=A(V\oplus V^*)\psi=SV,
$$
in particular $(S^iV^*)\psi=S^{m(V)-i}V$. 
 So taking $i=m(V)$ and using that $B$ gives an isomorphism 
$V\cong V^*$ (by evaluating on the first argument), we get an isomorphism $\psi^{\otimes 2}\cong \mathbf 1$. 
For general $i$ we get a surjective map $a_i: S^iV^*\otimes \psi\to S^{m(V)-i}V$. Thus the length 
of $S^iV^*$ (or, equivalently, $S^iV$) is at least the length of $S^{m(V)-i}V$. Replacing $i$ with $m(V)-i$, we obtain that 
these lengths are equal, hence $a_i$ is an isomorphism. 

Now, we claim that $SV$ is a Frobenius algebra. Indeed, we may assume that $\C$ is tensor-generated by $V$, then it has moderate growth by Lemma \ref{lemo}. So there exists a fiber functor $F: \C\to \sVec$  in characteristic zero (\cite{De}), respectively $F: \C\to \Ver_p$  in characteristic $p$ (\cite{CEO}), and it suffices to check that the algebra $SF(V)$ is Frobenius. In characteristic zero, 
$F(V)=E\otimes \psi$ for a finite dimensional vector space $E$, so $SF(V)$ is manifestly a Frobenius algebra. In characteristic $p$ 
it suffices to check that $SL_j$, $j>1$ is Frobenius; but this follows by applying the Verlinde fiber functor
$\Ver_p(SL(j))\to\Ver_p$ to the tautological object $\k^j$, whose
nonzero symmetric powers are simple and whose complementary-degree
multiplication pairings are nonzero.

Thus the multiplication map $S^iV\otimes S^{m(V)-i}V\to S^{m(V)}V=\psi$ is a non-degenerate pairing defining an isomorphism $S^{m(V)-i}V\otimes \psi\cong (S^iV)^*$. This implies (i).

(ii) By Lemma \ref{lemo}, we may assume that $\C$ has moderate growth. Thus 
by \cite{De},\cite{CEO}, there is a fiber functor $F: \C\to \sVec$ in characteristic zero and $F: \C\to \Ver_p$ 
in characteristic $p$. It suffices to prove the statement after applying $F$, i.e., in $\sVec$ or $\Ver_p$. 

In characteristic zero, $F(V)$ is purely odd, so $SV$ becomes an ordinary exterior algebra: $m(V)=\dim F(V)$ and its top line has parity $m(V)$, whence $c_{\psi,\psi}=(-1)^{m(V)}$. So assume that the characteristic is $p>0$ and use the decomposition in the proof of Corollary~\ref{simplic2}(i). For the summands $V_j=L_j\oplus L_j$ with $j$ odd we have 
$m(V_j)=2m(L_j)$, $S^{m(V_j)}V_j\cong S^{m(L_j)}L_j\otimes S^{m(L_j)}L_j\cong\mathbf1$. 
So both sides of the desired identity are $+1$.
For the remaining summands, recall that for $V=L_{2k}$ we have $m(V)=p-2k$ so $(-1)^{m(V)}=-1$, and $\psi=L_{p-1}$ (\cite{EOV}) so $c_{\psi,\psi}=-1$, i.e., (ii) holds. Since the statement of (ii) is compatible with direct sums (if it holds for $V_1,V_2$ then it does for $V_1\oplus V_2$), the claim follows.
\end{proof}

\subsection{Queer algebras in symmetric tensor categories}
Let $\C$ be a symmetric tensor category with a fixed super-line $\psi$. Recall that 
for $V\in \C$ the {\it queer algebra} $Q(V)$ is the image of the embedding 
$$
\underline{\End}(V)\oplus \underline\Hom(V,V\otimes \psi)\hookrightarrow \underline{\End}(V\oplus V\otimes \psi)
$$
 given by the schematic formula 
$(A,C)\mapsto \begin{pmatrix} A & C\\ C& A\end{pmatrix}$ (using that a morphism $C:V\to V\otimes\psi$ also gives a morphism $V\otimes\psi\to V$, which we also denote by $C$).

If $V\ne 0$, the algebras $\underline{\End}(V)$ and $Q(V)$ are simple, since they are 
Morita equivalent to simple algebras $\mathbf{1}$ and $Q(\mathbf{1})$. 
Namely, the $\C$-module category of $\underline{\End}(V)$-modules 
is equivalent to $\C$ via the assignment $X\in \C\mapsto X\otimes V$, while 
the $\C$-module category of $Q(V)$-modules is equivalent to 
$\C\boxtimes_\sVec\Vec$, where $\sVec$ is generated by $\psi$. 
The category $\C\boxtimes_\sVec\Vec$ is the category of objects 
$X\in\C$ equipped with an isomorphism $\gamma: X\to X\otimes \psi$ 
such that 
\[
(\id_X\otimes\alpha)\circ(\gamma\otimes\id_\psi)\circ\gamma=\id_X,
\]
which is the same as $Q(\mathbf{1})$-modules 
in $\C$, and the equivalence of categories $\C\boxtimes_\sVec\Vec\to Q(V)$-mod
is given by $X\mapsto X\otimes V$.
Under this equivalence $\psi\subset Q(\mathbf{1})=\mathbf{1}\oplus\psi$ acts by $\gamma$, while tensoring with $V$ implements the Morita equivalence from $Q(\mathbf{1})$ to $Q(V)$.

It is easy to see that $Q(V)\cong Q(V\otimes \psi)$ and 
$$
Q(V\otimes W)\cong Q(V)\otimes \underline{\End}(W)\cong \underline{\End}(V)\otimes Q(W),
$$
$$
Q(V)\otimes Q(W)\cong\underline{\End}(V\otimes W\oplus V\otimes W\otimes \psi).
$$

\subsection{Clifford and Weyl algebras in the Verlinde category}
Recall that $\Ver_p\cong \Ver_p^+\boxtimes \sVec$, where $\sVec$ is generated by $\psi=L_{p-1}$. 

\begin{proposition} The only simple algebras in $\Ver_p$ are of the form $\underline{\End} V$ and $Q(V)$, 
$0\ne V\in \Ver_p$. 
\end{proposition} 

\begin{proof} By \cite{EO,Os} the only indecomposable semisimple (=exact) module categories over $\Ver_p$ are $\Ver_p$ and $\Ver_p^+$.
Indeed, the ADET graphs with Coxeter number $p$ are $A_{p-1}$ and $T_{\frac{p-1}{2}}$, giving $\Ver_p$ and $\Ver_p^+$, respectively.

If $A$ is a simple algebra in $\Ver_p$ then $A$ is exact (\cite{CSZ}), so the category ${\mathcal M}:=A$-mod is an indecomposable exact (i.e., semisimple) module category over $\Ver_p$. Indecomposability follows because a decomposition of $A$-$\mathrm{mod}$ would give a nontrivial product decomposition of the algebra object $A$. Thus $\M\cong \Ver_p$ or $\M\cong \Ver_p^+$. 

Let $V$ be the object in $\M$ whose image is the regular left $A$-module. By the standard reconstruction of an algebra from a module category and a generator (cf. \cite[Subsection~7.10]{EGNO}),
$A\cong \underline{\End}_{\M}(V)$
as algebra objects of $\Ver_p$. Thus in the two cases $A\cong\underline{\End}(V)=V\otimes V^*$ and $A\cong Q(V)$, respectively.
\end{proof} 

\begin{proposition} Let $S_k$ be the image of the spin representation of $\operatorname{Spin}(2k-1)$ under the fiber functor $\Ver_p({\rm Spin}(2k-1))\to \Ver_p$.
 For each $1\le k\le \frac{p-1}{2}$ the Weyl algebra $A(L_{p-2k+1})=A_\psi(L_{2k-1})$ is isomorphic to $Q(S_k)$.\footnote{Here if $k=1$ then ${\rm Spin}(2k-1)=1$ and $S_k=\mathbf 1$. So in this case the proposition states that $A_\psi(\mathbf 1)\cong Q(\mathbf 1)$.}
\end{proposition} 

\begin{proof} 
 Let $U$ be the vector representation of ${\rm Spin}(2k-1)$ and $S$ its spin representation, and let ${\rm Cl}(U)$ be the (classical) Clifford algebra of $U$. In the classical tensor category of rational $\operatorname{Spin}(2k-1)$-modules, Clifford multiplication gives an isomorphism of equivariant associative superalgebras
$
\Cl(U)\cong Q(S).
$
For $2k-1<p$ the objects $U$ and $S$ are tilting and their images in the Verlinde category $\Ver_p(\operatorname{Spin}(2k-1))$ are non-negligible.
The fiber functor
$\Ver_p({\rm Spin}(2k-1))\longrightarrow \Ver_p$ (restriction along the principal homomorphism $SL(2)\to\operatorname{Spin}(2k-1)$ followed by semisimplification) sends $U$ to $L_{2k-1}$
 and $S$ to $S_k$, and it preserves the Clifford multiplication and the queer-algebra construction. Hence the image of the classical isomorphism is precisely $A_\psi(L_{2k-1})\cong Q(S_k)$.
Finally, since $L_{p-1}=\psi$ and $L_{2k-1}\otimes\psi\cong L_{p-2k+1}$ in $\Ver_p$, Lemma \ref{leweyl1}(i) identifies $A_\psi(L_{2k-1})$ with the Weyl algebra $A(L_{p-2k+1})$.
\end{proof} 

\subsection{The symplectic Witt group} 

Let $\C$ be a Frobenius exact symmetric tensor category. Let $V\in \C$ be a symplectic object, and assume that $SV$ has finite length. In this case by Corollary \ref{simplic2}(iii), $A(V)$ is an Azumaya algebra in $\C$. 
Thus it defines a class $C_V$ in the Brauer group $\Br(\C)$. 

Now assume that $\C$ is finite. Since $A(V)$ is a simple, hence exact, algebra in $\C$, the category $\C_V:=A(V)$-mod is an indecomposable exact $\C$-module category. Thus by Corollary~\ref{simplic2}(iii) $\C_V$ defines an element $\C_V$ in the Picard group $\Pic(\C)$ (\cite{ENO}, Subsection 4.4). In fact, the canonical isomorphism $\Br(\C)\cong \Pic(\C)$ maps $C_V$ to $\C_V$.  

Let ${\mathcal S\mathcal W}(\C)$ be the group generated by the classes $[V]$ of symplectic objects $V$ in $\C$ with $SV$ of finite length under direct sum, with defining relation $[V]=0$ if $A(V)\cong \underline{\End}S$ for some $S\in \C$. This is an elementary abelian 2-group (an $\Bbb F_2$-vector space). We call ${\mathcal S\mathcal W}(\C)$ the {\it symplectic Witt group} of $\C$. 

Let $\Br_2(\C)$ be the 2-torsion subgroup of $\Br(\C)$. 

\begin{proposition}\label{symwitt0} The assignment $V\mapsto C_V$ defines a subgroup inclusion 
$$
\iota: {\mathcal S\mathcal W}(\C)\hookrightarrow \Br_2(\C).
$$ 
\end{proposition} 

\begin{proof}
By Lemma~\ref{leweyl1}(ii),
\[
 A(V\oplus W)\cong A(V)\otimes A(W),
\]
so $V\mapsto[A(V)\text{-mod}]$ is additive. If $A(V)\cong\underline{\End}(S)$, its Brauer class is trivial, hence the defining relations of ${\mathcal S\mathcal W}(\C)$ lie in the kernel. Moreover, since $A(V)\cong A(V)^{\rm op}$, 
\[
 A(V)\otimes A(V)\cong A(V)\otimes A(V)^{\rm op}
 \cong\underline{\End}(A(V)),
\]
where the second isomorphism is the Azumaya property from Corollary~\ref{simplic2}(iii); thus $2[V]$ maps to zero. This gives a homomorphism 
$$
\iota:{\mathcal S\mathcal W}(\C)\to\Br_2(\C).
$$

It remains to check injectivity. If a finite sum $[V_1]+\cdots+[V_r]$ maps to zero, set $V=V_1\oplus\cdots\oplus V_r$. Then $A(V)$ is Brauer-trivial. An equivalence $A(V)\text{-mod}\simeq\C$ sends the free rank-one module to an object $S\in\C$, and reconstruction gives $A(V)\cong\underline{\End}(S)$. Hence $[V]=0$ is already one of the defining relations of ${\mathcal S\mathcal W}(\C)$. Therefore $\iota$ is injective.
\end{proof}

\subsection{The symplectic Witt group of ${\rm Rep}(G)\boxtimes \sVec$ for a finite group $G$} 
Let $G$ be a finite group whose order is coprime to ${\rm char}(\k)$ if it is $>0$. Let $\C={\rm Rep}(G)\boxtimes \sVec$ and $\Gamma:=G\times \Bbb Z/2$ where $\Bbb Z/2$ is generated by an element $t$.
 Thus $\C={\rm Rep}(\Gamma,t)$, the super-Tannakian category of representations of $G$ on supervector spaces; the element $t$ is defined to act on such a representation by the parity operator. 
In this case by Carnovale's theorem (\cite{DN}, Theorem 6.5) 
$$
\Pic(\C)=H^2_*(\Gamma,\k^\times)\times \mu_2,
$$ 
with $H^2_*$ denoting $H^2$ with modified multiplication
 defined at the level of 2-cocycles 
by the formula
$$
(\beta*\gamma)(x,y)=(-1)^{\varepsilon_\beta(x)\varepsilon_\gamma(y)}\beta(x,y)\gamma(x,y),
$$
where $\varepsilon_\beta: \Gamma\to \Bbb Z/2$  is the homomorphism 
 such that
 $$
 (-1)^{\varepsilon_\beta(x)}=\beta(x,t)/\beta(t,x)\in \mu_2
$$
(it depends only on the cohomology class of $\beta$). 

If $X$ is a representation of $G$, denote by $X_-$ its extension to a representation of $\Gamma$ by letting $t$ act by $-1$. 
Also recall that ${\mathcal S\mathcal W}(\C)$ is generated by $[X\otimes \psi]$ 
where $X$ is an orthogonal representation of $G$. Let 
$\widetilde X:=X_-\oplus \k_-^{\dim X}$; thus the action of $\Gamma$ on $\widetilde X$ is given by 
$$
\widetilde\rho(g,t^a):=(-1)^a(\rho(g)\oplus \id_{\k^{\dim X}})\in O(\widetilde X).
$$ 
This homomorphism gives rise to a class $\beta_X\in H^2_*(G\times \Bbb Z/2,\k^\times)$ 
which is the pullback by $\widetilde\rho$ of the central extension 
$$
1\to \k^\times\to {\rm Cliff}(\widetilde X)\to O(\widetilde X)\to 1,
$$
where ${\rm Cliff}(\widetilde X)$ is the Clifford group of $\widetilde X$. 

\begin{proposition}\label{symwitt} The map $\iota$ is given by 
$$
\iota([X\otimes \psi])=(\beta_X,(-1)^{\dim X}),
$$ 
and it is a group homomorphism. Thus $\iota({\mathcal S\mathcal W}(\C))=\mathcal S\mathcal W_+(\C)\times \mu_2$, where $\mathcal S\mathcal W_+(\C)$ 
is the subgroup of $H^2_*(G\times \Bbb Z/2,\k^\times)$ generated by the $\beta_X$. 
\end{proposition} 

\begin{proof} Under the equivalence
\[
\operatorname{Rep}(G)\boxtimes\sVec\simeq \operatorname{Rep}(\Gamma,t),
\]
the algebra $A(X\otimes\psi)=A_\psi(X)$ is the usual Clifford superalgebra $\Cl(X_-)$. 
Write $n=\dim X$. The image of $[\Cl(X_-)\text{-mod}]$ under restriction to $\operatorname{Pic}(\sVec)=\mu_2$ is $(-1)^n$: after forgetting the $G$-action, $\Cl(X_-)$ is the graded tensor product of $n$ copies of the one-generator odd Clifford algebra, which represents the nontrivial element of $\operatorname{Pic}(\sVec)$.

We next compute the first coordinate. Choose homogeneous lifts $u_\gamma\in\operatorname{Cliff}({\widetilde X})$ of $\widetilde\rho(\gamma)$, $\gamma\in\Gamma$, and write
\[
u_\gamma u_\delta=b(\gamma,\delta)u_{\gamma\delta}.
\]
Then $[b]=\beta_X$. Let $p(\gamma)=|u_\gamma|\in \Bbb Z/2$, i.e., $\det \widetilde \rho(\gamma)=(-1)^{p(\gamma)}$; thus $p: \Gamma\to \Bbb Z/2$ is a group homomorphism. Choose a normalized volume element $\omega\in\Cl({\widetilde X})$ with $\omega^2=1$. It is even and satisfies
\[
\omega a=(-1)^{|a|}a\omega
\]
for every homogeneous $a\in\Cl({\widetilde X})$. Choose $i\in\k$ with $i^2=-1$ and define
\[
T_\gamma:=i^{p(\gamma)}\omega^{p(\gamma)}u_\gamma.
\]
The Clifford-group map is the twisted adjoint map, so for $v\in {\widetilde X}$,
\[
\widetilde\rho(\gamma)(v)=(-1)^{p(\gamma)}u_\gamma vu_\gamma^{-1}.
\]
It follows from the defining property of $\omega$ that
\[
T_\gamma vT_\gamma^{-1}=\widetilde\rho(\gamma)(v).
\]
Moreover, since $u_\gamma\omega^{p(\delta)}=(-1)^{p(\gamma)p(\delta)}\omega^{p(\delta)}u_\gamma$, $\omega^2=1$ and $i^2=-1$, we obtain
\[
T_\gamma T_\delta=b(\gamma,\delta)T_{\gamma\delta}.
\]
Now, since $\widetilde X$ is even-dimensional, $\Cl({\widetilde X})$ is a matrix superalgebra, therefore $\Cl({\widetilde X})\cong\underline{\operatorname{End}}(S)$ for a graded spinor module $S$. Since $\widetilde\rho(t)=-\id_{\widetilde X}$ and $p(t)=0$, we may normalize the lift $u_t$ to be the volume element $\omega$; then $T_t=\omega$ acts on $S$ as the parity operator. Thus the grading on $S$ is exactly the one determined by $t$, as required in the projective-representation description of the $H^2$-coordinate in the proof of \cite[Theorem~6.5]{DN}.
 Under this isomorphism the $\Gamma$-action on $\Cl({\widetilde X})$ is induced by the projective action $T_\gamma$ on $S$, whose multiplier is $b$. By the kernel description in the proof of \cite[Theorem~6.5]{DN}, it follows that
\[
[\Cl({\widetilde X})\text{-mod}]=(\beta_X,1).
\]

On the other hand, the orthogonal direct-sum formula for Clifford algebras gives a $\Gamma$-equivariant isomorphism
\[
\Cl({\widetilde X})\cong\Cl(X_-)\otimes\Cl(\k_-)^{\otimes n},
\]
where the tensor products are graded. The algebra $\Cl(\k_-)=\mathbf 1\oplus\k_-$ represents $(1,-1)$: its restriction to $\sVec$ is the nontrivial Brauer--Wall class, while its first coordinate is trivial because its $\Gamma$-action factors through $\langle t\rangle$ and $H^2(\langle t\rangle,\k^\times)$ vanishes. Consequently
\[
(\beta_X,1)=[\Cl({\widetilde X})\text{-mod}]
=[\Cl(X_-)\text{-mod}]\,(1,-1)^n,
\]
and hence\footnote{Since $\iota$ is a group homomorphism by Proposition~\ref{symwitt0},
this formula is compatible with direct sums.}
\[
\iota([X\otimes\psi])=[\Cl(X_-)\text{-mod}]
=(\beta_X,(-1)^n).
\]

Finally, for $X=\k$, the stabilized representation ${\widetilde X}=\k_-\oplus\k_-$ factors through $\langle t\rangle$, so $\beta_\k=1$ and
\[
\iota([\psi])=(1,-1).
\]
Thus the full factor $\mu_2$ lies in the image. For arbitrary $X$,
\[
(\beta_X,1)=\iota([X\otimes\psi])\,(1,-1)^n
\]
also lies in the image. Hence, identifying ${\mathcal S\mathcal W}(\C)$ with its image under $\iota$, that image is exactly $\mathcal S\mathcal W_+(\C)\times\mu_2$.
\end{proof}

\begin{remark}\label{precrel} The proof of Proposition \ref{symwitt} gives the precise relation
 with the pullback of the Clifford-group extension for the unstabilized representation $X_-$.  Namely, the direct-product decomposition $\Gamma=G\times\langle t\rangle$ gives an isomorphism of ordinary cohomology groups
\[
\Theta:H^2(\Gamma,\k^\times)\xrightarrow{\ \sim\ }
H^2(G,\k^\times)\times\Hom(G,\mu_2)
\]
characterized by
\[
\Theta([c])=
\left([c|_{G\times G}],\ g\longmapsto\frac{c(t,g)}{c(g,t)}\right).
\]
Indeed, $H^2(\langle t\rangle,\k^\times)$ vanishes, and the second factor is the mixed term.  Let $b_X^{\rm raw}$ be the pullback class along
\[
\Gamma\longrightarrow O(X),\qquad (g,t^a)\longmapsto(-1)^a\rho(g),
\]
and let $b_G(X)\in H^2(G,\k^\times)$ be the pullback class along $\rho$.  More generally, if $U$ is an $m$-dimensional orthogonal $\Gamma$-representation on which $t$ acts by $-1$, a volume element $\omega_U$ is a Clifford lift of $t$, and for a homogeneous lift $u_g$ of $g\in G$ one has
\[
\omega_Uu_g=(-1)^{(m-1)|u_g|}u_g\omega_U,
\qquad
(-1)^{|u_g|}=\det(g|_U).
\]
Consequently
\[
\frac{\widetilde b_U(t,g)}{\widetilde b_U(g,t)}=\det(g|_U)^{m-1},
\]
where $\widetilde b_U$ is the cocycle obtained from the chosen Clifford lifts for $b_U^{\rm raw}$.
Applying this first to $U=X_-$ and then to $U=\widetilde X$, and observing that the stabilizing summand is $G$-trivial, gives
\[
\Theta(b_X^{\rm raw})=\bigl(b_G(X),{\rm det}_X^{n-1}\bigr),
\qquad
\Theta(\beta_X)=\bigl(b_G(X),{\rm det}_X\bigr).
\]
For $\chi\in\Hom(G,\mu_2)$ define the mixed cocycle
\[
\kappa_\chi\bigl((g,t^a),(h,t^b)\bigr):=\chi(h)^a.
\]
It has trivial restriction to $G\times G$ and satisfies
\[
\frac{\kappa_\chi(t,g)}{\kappa_\chi(g,t)}=\chi(g).
\]
Thus
\[
\beta_X=b_X^{\rm raw}\,\kappa_{\det_X}^{\dim X}
\quad\text{in }H^2(\Gamma,\k^\times).
\]
In particular, the raw pullback is already the first Picard coordinate when $n$ is even, whereas in odd dimension it must be corrected by $\kappa_{\det_X}$.

For two orthogonal $G$-representations $X,Y$, the isomorphism
\[
\Cl((X\oplus Y)_-)\cong\Cl(X_-)\otimes\Cl(Y_-)
\]
shows that these classes multiply. Equivalently,
\[
\beta_{X\oplus Y}=\beta_X*\beta_Y,
\]
and the second coordinates satisfy
\[
(-1)^{\dim(X\oplus Y)}=(-1)^{\dim X}(-1)^{\dim Y}.
\]
This gives a direct proof that the formula of Proposition \ref{symwitt} defines a group homomorphism.
\end{remark} 

\subsection{The Stiefel--Whitney subgroup}

Let $\k=\Bbb C$. Put
\[
A:=H^2(G,\k^\times),\qquad K:=\Hom(G,\Bbb Z/2)=H^1(G,\Bbb Z/2),
\]
and let
\[
j:H^2(G,\Bbb Z/2)\longrightarrow H^2(G,\k^\times)
\]
be induced by the inclusion $\Bbb Z/2\hookrightarrow\k^\times$. Under the ordinary K\"unneth isomorphism
\[
\Theta:H^2(\Gamma,\k^\times)\xrightarrow{\sim}A\times K
\]
from Remark \ref{precrel}, the modified multiplication is
\begin{equation}\label{modifiedKunneth}
(a,\chi)*(b,\eta)=\bigl(abj(\chi\smile\eta),\chi+\eta\bigr).
\end{equation}
Indeed, the extra sign in the definition of $*$ restricts to $G\times G$ as the cocycle representing $j(\chi\smile\eta)$. 
Thus, we have a canonical short exact sequence of abelian groups 
$$
1\to H^2(G,\k^\times)\to H^2_*(G\times \Bbb Z/2,\k^\times)\to \Hom(G,\Bbb Z/2)\to 1.
$$
Notice that
\begin{equation}\label{cupsquarezero}
j(\chi\smile\chi)=1
\end{equation}
for every $\chi\in K$; indeed, the cocycle $(-1)^{\chi(g)\chi(h)}$ is the coboundary of the cochain $g\mapsto i^{\chi(g)}$, where $i^2=-1$.

For an orthogonal representation $X$ of $G$, let $w_1(X)\in K$ and $w_2(X)\in H^2(G,\Bbb Z/2)$ be its first and second Stiefel--Whitney classes of $X$ (i.e., of the underlying orthogonal vector bundle $E_X$ on $BG$); thus $\det_X=(-1)^{w_1(X)}$. Then
\[
b_G(X):=j(w_2(X))\in A[2],
\qquad
\Theta(\beta_X)=\bigl(b_G(X),w_1(X)\bigr).
\]
Define
\[
{\mathcal S\mathcal W}(G):=
\left\langle b_G(X)\ \middle|\ X\text{ is an orthogonal representation of }G\right\rangle
\subset A[2].
\]
We call this the \emph{Stiefel--Whitney subgroup} of $G$.

\begin{proposition}\label{SWproductstructure} (Symplectic Witt via Stiefel-Whitney)
Under $\Theta$ one has
\[
\mathcal S\mathcal W_+(\C)={\mathcal S\mathcal W}(G)\times K
\]
as a set. In particular, $\mathcal S\mathcal W_+(\C)$ is also a subgroup of $H^2(\Gamma,\k^\times)$ for the ordinary multiplication.

Moreover, let
\[
D(G):=
\left\langle j(\chi\smile\eta)\ \middle|\ \chi,\eta\in K\right\rangle
\subset A.
\]
Then $D(G)\subset {\mathcal S\mathcal W}(G)$, and
\[
L:=D(G)\times K\subset \mathcal S\mathcal W_+(\C)
\]
is a subgroup for the modified multiplication. Namely, it is the subgroup generated by the elements $(1,\chi)$, $\chi\in K$. 
\end{proposition}

\begin{proof} The Whitney sum formula
$$
w_2(X\oplus Y)=w_2(X)+w_2(Y)+w_1(X)\smile w_1(Y)
$$
yields
\begin{equation}\label{Whitneysumdegree2}
b_G(X\oplus Y)=b_G(X)b_G(Y)j\bigl(w_1(X)\smile w_1(Y)\bigr).
\end{equation}
Since every element of $K$ is $w_1$ of an orthogonal line, it follows that $D(G)\subset {\mathcal S\mathcal W}(G)$. Formula \eqref{modifiedKunneth} and Proposition~\ref{symwitt} therefore imply
\[
\mathcal S\mathcal W_+(\C)\subset {\mathcal S\mathcal W}(G)\times K.
\]
Conversely, for every $\chi\in K$, the corresponding orthogonal line $\k_\chi$ satisfies
\[
\Theta(\beta_{\k_\chi})=(1,\chi).
\]
Hence, by \eqref{modifiedKunneth} and \eqref{cupsquarezero},
\[
\Theta(\beta_X)*(1,w_1(X))=(b_G(X),0).
\]
Thus $\mathcal S\mathcal W_+(\C)$ contains ${\mathcal S\mathcal W}(G)\times\{0\}$ and $\{1\}\times K$, and hence contains their Cartesian product. This proves the asserted equality of sets, and the statement about the ordinary multiplication follows immediately.

The subset $L$ is closed under $*$ by \eqref{modifiedKunneth}, and every element of $L$ is its own inverse by \eqref{cupsquarezero}; hence it is a subgroup. Also,
\[
(1,\chi)*(1,\eta)*(1,\chi+\eta)=\bigl(j(\chi\smile\eta),0\bigr),
\]
so $L$ is generated by the elements $(1,\chi)$. 
\end{proof}

\begin{remark} Note that the subset $K=\{(1,\chi)\mid\chi\in K\}$ need not be a subgroup for $*$. Indeed, 
\[
(1,\chi)*(1,\eta)=\bigl(j(\chi\smile\eta),\chi+\eta\bigr),
\]
which need not belong to $K$. For example, for $G=(\Bbb Z/2)^2$ and the two coordinate characters $\chi,\eta$, the class $j(\chi\smile\eta)$ is nonzero, as is seen from the nontrivial commutator of its representing cocycle.
\end{remark} 


\begin{corollary}\label{SWsurjectivitycriterion}
(i) One has, as sets,
\[
H^2_*(\Gamma,\k^\times)[2]=H^2(G,\k^\times)[2]\times K
\]
and
\[
\iota\bigl({\mathcal S\mathcal W}(\C)\bigr)
={\mathcal S\mathcal W}(G)\times K\times\mu_2.
\]
Consequently,
\[
\iota:{\mathcal S\mathcal W}(\C)\longrightarrow\Br_2(\C)
\]
is surjective if and only if
\[
{\mathcal S\mathcal W}(G)=H^2(G,\k^\times)[2].
\]

(ii) $\iota$ is surjective if $G$ is abelian. 

(iii) $\iota$ is not always surjective.
\end{corollary}

\begin{proof}
(i) By \eqref{modifiedKunneth} and \eqref{cupsquarezero},
\[
(a,\chi)*(a,\chi)=(a^2,0),
\]
which proves the first equality. The second follows from Proposition~\ref{symwitt} and Proposition~\ref{SWproductstructure}, and gives the desired criterion.

(ii) If $G$ is abelian then $H^2(G,\k^\times)[2]=\wedge^2 H^1(G,\Bbb Z/2)$, so 
the map 
$$
j\circ \smile: H^1(G,\Bbb Z/2)\otimes H^1(G,\Bbb Z/2)\to H^2(G,\k^\times)[2]
$$ 
is surjective. Hence $L=A[2]\times K$ and thus $\iota$ is surjective. 

(iii) The type-$1$ split metacyclic groups occurring in the proof of \cite[Theorem~2, pp.~336--337]{GKT} have a degree-two class outside the subgroup generated by second Stiefel--Whitney classes. Its image under
\[
j:H^2(G,\Bbb Z/2)\longrightarrow H^2(G,\k^\times)[2]
\]
lies outside ${\mathcal S\mathcal W}(G)$. Indeed, if $j(x)$ belonged to ${\mathcal S\mathcal W}(G)$, then $x$ would differ from a sum of second Stiefel--Whitney classes by an element of $\ker j$; but $\ker j$ consists of Kummer classes of complex characters, which are themselves second Stiefel--Whitney classes of the underlying oriented real two-planes. Thus ${\mathcal S\mathcal W}(G)\subsetneq H^2(G,\mathbb C^\times)[2]$, and the result follows from (i).
\end{proof}

\begin{thebibliography}{999999}

\bibitem[CEO]{CEO} K.~Coulembier, P.~Etingof, and V.~Ostrik, \emph{On Frobenius exact symmetric tensor categories}, with an appendix by A.~Kleshchev, Ann. of Math. (2) \textbf{197} (2023), no.~3, 1235--1279; doi:10.4007/annals.2023.197.3.5; arXiv:2107.02372.

\bibitem[CEN]{CEN} K.~Coulembier, P.~Etingof, and J.~Newton, \emph{Finite symmetric algebras in tensor categories and Verlinde categories of algebraic groups}, Adv. Math. \textbf{483} (2025), 110677; doi:10.1016/j.aim.2025.110677; arXiv:2502.10598.

\bibitem[CSZ]{CSZ}  K.~Coulembier, M.~Stroi\'{n}ski, and T.~Zorman, \emph{Simple algebras and exact module categories}, arXiv:2501.06629v2 [math.RT] (2025).

\bibitem[DN]{DN} A.~Davydov and D.~Nikshych, \emph{Braided Picard groups and graded extensions of braided tensor categories}, Selecta Math. (N.S.) \textbf{27} (2021), no.~4, Paper No.~65, 87~pp.; doi:10.1007/s00029-021-00670-1; arXiv:2006.08022.

\bibitem[De]{De} P.~Deligne, \emph{Cat\'egories tensorielles}, Mosc. Math. J. \textbf{2} (2002), no.~2, 227--248; doi:10.17323/1609-4514-2002-2-2-227-248.

\bibitem[E]{E} P.~Etingof, \emph{Koszul duality and the PBW theorem in symmetric tensor categories in positive characteristic}, 
\emph{Adv. Math.} \textbf{327} (2018), 128--160; doi:10.1016/j.aim.2017.06.014; arXiv:1603.08133.

\bibitem[EGNO]{EGNO} P.~Etingof, S.~Gelaki, D.~Nikshych, and V.~Ostrik, \emph{Tensor categories}, Mathematical Surveys and Monographs, vol.~205, American Mathematical Society, Providence, RI, 2015; doi:10.1090/surv/205.

\bibitem[ENO]{ENO} P.~Etingof, D.~Nikshych, and V.~Ostrik, with an appendix by E.~Meir, \emph{Fusion categories and homotopy theory}, Quantum Topol. \textbf{1} (2010), no.~3, 209--273; doi:10.4171/QT/6; arXiv:0909.3140.

\bibitem[EO]{EO} P.~Etingof and V.~Ostrik, \emph{Module categories over representations of $SL_q(2)$ and graphs}, Math. Res. Lett. \textbf{11} (2004), no.~1, 103--114; doi:10.4310/MRL.2004.v11.n1.a10; arXiv:math/0302130.

\bibitem[EOV]{EOV} P.~Etingof, V.~Ostrik, and S.~Venkatesh, \emph{Computations in symmetric fusion categories in characteristic $p$}, Int. Math. Res. Not. IMRN \textbf{2017} (2017), no.~2, 468--489; doi:10.1093/imrn/rnw024; arXiv:1512.02309.

\bibitem[GKT]{GKT} J.~Gunarwardena, B.~Kahn, and C.~B.~Thomas, \emph{Stiefel--Whitney classes of real representations of finite groups}, J. Algebra \textbf{126} (1989), no.~2, 327--347; doi:10.1016/0021-8693(89)90309-8.

\bibitem[Os]{Os} V.~Ostrik, \emph{Module categories over representations of $SL_q(2)$ in the non-semisimple case}, Geom. Funct. Anal. \textbf{17} (2008), no.~6, 2005--2017; doi:10.1007/s00039-007-0637-4; arXiv:math/0509530.

\bibitem[OZ]{OZ} F.~Van Oystaeyen and Y.~H.~Zhang, \emph{The Brauer group of a braided tensor category}, J. Algebra \textbf{202} (1998), no.~1, 96--128; doi:10.1006/jabr.1997.7295.



\end{thebibliography}

\end{document}

